% 负数（综述）
% license CCBYSA3
% type Wiki

本文根据 CC-BY-SA 协议转载翻译自维基百科\href{https://en.wikipedia.org/wiki/Negative_number}{相关文章}

\begin{figure}[ht]
\centering
\includegraphics[width=6cm]{./figures/18d7e2ba5a51f4aa.png}
\caption{} \label{fig_FS1_1}
\end{figure}
在数学中，负数是正实数的相反数。[1]等价地说，负数是小于零的实数。负数常用于表示损失或缺少的量。例如，所欠的债务可以看作一种负资产。如果某个量（如电子的电荷）可以具有两种相反的性质，那么人们可以（有时是任意地）将其中一种记为正，另一种记为负。负数还用于描述那些可以小于零的刻度值，例如摄氏温度和华氏温度的刻度。负数的算术运算规律保证了“相反数”的直观概念能够反映在运算中。例如：$-(-3) = 3$因为“相反数的相反数”就是原数本身。

负数通常写作前面带一个负号。例如：$-3$ 表示一个大小为 3 的负量，读作“负三”或“减三”。相反地，大于零的数称为正数；零通常（但并非总是）被认为既不是正数，也不是负数。[2]正数的性质有时会用加号来强调，例如：$+3$。一般而言，数的“正负性”被称为它的符号。


除了零之外，每一个实数要么是正数，要么是负数。非负整数通常称为自然数（即 $0, 1, 2, 3, \ldots$），而正整数、负整数以及零合在一起称为整数。（有些自然数的定义并不包括零。）

在簿记中，所欠的金额常用红色数字或括号中的数字表示，作为负数的一种替代记号。负数早在中国汉代（公元前202年 – 公元220年）成书的《九章算术》中就已使用，不过其中可能包含更古老的材料。[3]大约三世纪的数学家刘徽建立了负数加减的运算规则。[4]到七世纪，印度数学家如婆罗摩笈多已经在描述负数的运用。伊斯兰数学家进一步发展了负数的减法与乘法规则，并解决了带有负系数的问题。[5]在负数概念出现之前，像丢番图这样的数学家认为负解是“错误的”，并称要求负解的方程为荒谬的。[6]西方数学家如莱布尼茨曾坚持认为负数是无效的，但在计算中仍然使用它们。[7][8]

